\documentclass[10pt,a4paper]{amsart}
\usepackage[margin=19mm]{geometry}
\usepackage{amsmath,amssymb,mathtools}
\usepackage[scr=rsfs,cal=euler]{mathalfa}
\usepackage{xfrac}
\usepackage[unicode,hidelinks,bookmarks=false]{hyperref}
\newcommand{\ie}{i.e.\ }
\newcommand{\cf}{cf.\ }
\newcommand{\resp}{respectively\ }
\newcommand{\Subsection}{\subsection}
\providecommand{\nobreakdash}{\nobreak\hbox{-}}
\setlength{\emergencystretch}{2em}
\multlinegap0pt
\numberwithin{equation}{section}
\newtheorem{theorem}{Theorem}[section]
\newtheorem{lemma}[theorem]{Lemma}
\newtheorem{remark}[theorem]{Remark}
\allowdisplaybreaks
\title[Two nontrivial solutions for a fractional p-Laplacian Kirchhoff]{Two nontrivial solutions for a fractional p-Laplacian Kirchhoff problem with steep potential well}
\author{Shuwen He, Shiqing Zhang}
\date{Source: Electronic Journal of Differential Equations, Vol. 2026 (2026), No. 15, pp. 1--19; DOI 10.58997/ejde.2026.15; published February 17, 2026; publisher version}
\begin{document}
\maketitle
\noindent\emph{Electronic Journal of Differential Equations}, Vol. 2026 (2026), No. 15, pp. 1--19.\newline
ISSN: 1072-6691. DOI \href{https://doi.org/10.58997/ejde.2026.15}{10.58997/ejde.2026.15}. This work is licensed under a CC BY 4.0 license.\par\smallskip
\noindent\textit{Editorial scope.} Counted claim: original Theorem 1.1 of the published article (native label thm1.1, source lines 238--243, printed as Theorem 1.1 in the published PDF), namely that under (A1)--(A8) and for each fixed a there are thresholds b-hat, Lambda and Pi such that the fractional p-Laplacian Kirchhoff problem has at least two different nontrivial solutions, together with its complete original author proof printed as Proof of Theorem 1.1 at source lines 1132--1189. Necessary complete context is retained uncounted: the whole Introduction with hypotheses (A1)--(A8) and the statements of Theorems 1.1--1.5 and Remark 1.6 (source lines 61--237), and the whole of Section 2 on preliminaries and Section 3 on existence, including the complete proofs of Lemmas 3.1--3.7 and Remark 3.4 that the counted proof invokes (source lines 245--1131). Nothing inside the delivered ranges is omitted, so every printed section, theorem, lemma, remark and equation number of the delivered unit coincides with the published PDF. The delivered unit ends after the counted proof with the reference list of the cited published entries. Printed slips kept literal: the statement of Theorem 1.2 prints the clause $b\in(0,\widehat{b})$ together with the convergence of the two solutions and then has no second clause after the word and; the statements of Theorems 1.3 and 1.4 print the superscript of the two solutions as a plain $(i=1,2)$ attached to the symbol $u_{b,\lambda}$ while the body of the article prints the parenthesised superscript form $u^{(i)}_{b,\lambda}$; and the chain $0<\eta\leq c_{b,\lambda}^{\beta}\leq\widehat{c}$ is printed twice, once inline at source line 1125 without the strict upper bound and once as the numbered display (3.24) at source lines 1127--1130. All of those printed forms are delivered unchanged. No mathematics was written, repaired or re-derived, and no figure, image, table or external asset occurs anywhere in the delivered ranges.
\section{Introduction}

The aim of this article is to study the fractional $p$-Laplacian Kirchhoff type 
problems with steep potential well and concave-convex nonlinearities,
\begin{equation}\label{lab1.1}
\begin{gathered}
\Big(a+b\int\int_{\mathbb{R}^{2N}}\frac{|u(x)-u(y)|^p}{|x-y|^{N+sp}}\,dx\,dy\Big)
(-\Delta)_p^su+\lambda V(x)|u|^{p-2}u=f(x,u)+g(x,u)\quad \text{in }\mathbb{R}^N, \\
u\in W^{s,p}(\mathbb{R}^N),
\end{gathered}
\end{equation}
where $s\in(0,1), 2\leq p<\infty, N>sp, a, b, \lambda>0$ are real parameters, $V, f$ 
and $g$ are continuous functions, $(-\Delta)^s_p$ is the fractional $p$-Laplacian 
which (up to normalization factors) is the nonlinear nonlocal operator defined on 
smooth functions by
$$
(-\Delta)^s_pu(x)=2\lim_{\epsilon\to 0}\int_{\mathbb{R}^N\setminus 
B_\epsilon(x)}\frac{|u(x)-u(y)|^{p-2}[u(x)-u(y)]}{|x-y|^{N+sp}}dy\quad
 \text{in }\mathbb{R}^N,
$$
where $u\in C_0^\infty(\mathbb{R}^N)$, 
$B_\epsilon(x):=\{y\in\mathbb{R}^N:|y-x|<\epsilon\}$. Nonlocal fractional operators 
come from many different contexts, such as finance, quantum mechanics, 
game theory and so on, see for example \cite{r5, r14, r15,r18} and the references 
therein. One of the main reasons of considering \eqref{lab1.1} is related to the 
equation
$$
 \rho u_{tt}-\Big(\frac{p_0}{h}+\frac{E}{2L}\int^L_0 |u_x|^2dx \Big)u_{xx}=w(x,u),
$$
which was proposed by Kirchhoff in \cite{r13} as an extension of the classical 
D'Alembert wave equations for free vibrations of elastic strings. 
This model takes into account the changes in length of the string produced by 
transverse vibrations. For more mathematical and physical background on the 
Kirchhoff type problems, we refer the readers to \cite{r1, r4} and the references 
therein.

In recent years, when parameters $a, b, \lambda>0$, $V$ and $f$ satisfy some 
appropriate conditions, the general Kirchhoff type problem
\begin{gather*}
-\Big(a+b\int_{\mathbb{R}^N}|\nabla u|^2dx\Big)\Delta u+\lambda V(x)u= f(x,u)\quad
 \text{in } \mathbb{R}^N, \\
u\in H^1(\mathbb{R}^N)
\end{gather*}
has been widely studied by many researchers, they obtained the existence, 
concentration, and multiplicity of solutions to this equation through variational 
methods and analytical techniques see \cite{r10, r16, r23, r27, r30, r34} and the 
references therein.

More recently, Xiang et al.\  \cite{r31} investigated the  nonlocal fractional 
$p$-Laplacian Kirchhoff type problem
\begin{equation}\label{lab1.2}
\begin{gathered}
M\Big(\int\int_{\mathbb{R}^{2N}}|u(x)-u(y)|^p K(x-y)\,dx\,dy\Big)
\mathcal{L}_K^pu=f(x,u)\quad \text{in }\Omega, \\
u=0\quad \text{in }\mathbb{R}^N\setminus\Omega,
\end{gathered}
\end{equation}
where $M$ is a continuous function, $\mathcal{L}_K^p$ is a nonlocal operator with 
singular kernel $K$. When Carath\'{e}odory function $f$ is sublinear growth 
or superlinear growth, the authors obtained the existence nontrivial solutions 
separately. Additionally, when $p=2$, Fiscella and Valdinoci  \cite{r8} established 
the existence of nontrivial solutions to problem \eqref{lab1.2} with a critical 
growth term. In particular, a typical example for $K$ is given by singular kernel 
$K(x-y)=|x-y|^{-(N+sp)}$ and $V(x)|u|^{p-2}u, g(x)$ are added to \eqref{lab1.2} 
defined on the unbounded domain $\mathbb{R}^N$. In this case, problem \eqref{lab1.2} 
becomes
\begin{equation}\label{lab1.3}
M\Big(\int\int_{\mathbb{R}^{2N}}\frac{|u(x)-u(y)|^p}{|x-y|^{N+sp}}\,dx\,dy\Big)
(-\Delta)_p^s u+V(x)|u|^{p-2}u=f(x,u)+g(x)\quad \text{in }\mathbb{R}^N.
\end{equation}
If $g(x)\in L^{\frac{p}{p-1}}(\mathbb{R}^N)$ is a nonzero perturbation term and 
$V$ and $f$ satisfy the following assumptions
\begin{itemize}
\item[(1)] $\inf_{x\in\mathbb{R}^N}V(x)\geq V_0>0$ for the positive constant $V_0$;

\item[(2)] there exists $R>0$ such that $\lim_{|y|\to \infty}|
\{x\in B_{R}(y): V(x)\leq d \} |=0$ for every $d>0$,
where $|\cdot|$ is the Lebesgue measure in $\mathbb{R}^N$;

\item[(3)] there is $\mu>p$ such that
\[
 \mu \int_0^tf(x,\tau)d\tau\leq f(x,t)t \quad \text{for all }  
 (x, t)\in \mathbb{R}^N\times\mathbb{R}.
\]
\end{itemize}
Pucci et al.\  \cite{r19} proved that problem \eqref{lab1.3} admits at least two 
nontrivial solutions. Subsequently, Torres Ledesma  \cite{r26} also obtained the 
existence of two nontrivial solutions to problem \eqref{lab1.3} when $f$ satisfies
\begin{itemize}
\item[(4)] there exist $q\in(p, p^*_s)$ and $h(x)\in C(\mathbb{R}^N, \mathbb{R}^+)$ 
with $\lim_{|x|\to \infty}h(x)=0$ such that
\[
 |f(x,t)|\leq h(x)|t|^{q-1} \quad \text{for all } 
  (x, t)\in \mathbb{R}^N\times\mathbb{R}.
\]
\end{itemize}

Recently, when $f(x,u)+g(x)=\lambda\omega(x)|u|^{r-2}u-h(x)|u|^{q-2}u$, where 
$h(x)$ is a nonnegative function satisfying some ratio of integration with
$\omega(x)$, $1 < r < q <\infty$, Pucci et al.\ showed the existence and multiplicity 
of entire solutions for problem \eqref{lab1.3} in \cite{r20}. 
Later, for the critical case of $f(x,u)+g(x)=\alpha|u|^{p^*_s-2}u+\beta 
\kappa(x)|u|^{q-2}u$, Wang and Zhang  \cite{r28} showed the existence of infinitely 
many solutions which tend to zero under the appropriate positive parameters $\alpha$ 
and $\beta$. For more results on different forms of problem \eqref{lab1.3}, readers 
can see  \cite{r9, r17, r25, r32, r35} and the references therein. However, 
there seem to be few articles studying fractional $p$-Laplacian Kirchhoff type problems 
with steep potential well and concave-convex nonlinearities. We note that, when the 
concave-convex term presents as $f(x)|u|^{q-2}u+g(x)|u|^{r-2}u$ with $1<r<p<q<p^*_s$, 
Xiong et al.\  \cite{r33} explored the existence and multiplicity of nontrivial 
solutions to problem \eqref{lab1.1}. Specifically, if $p=2$, $g(x,u)=\mu g(x)u^q$ for all 
$\lambda>0$ and $0<q<1$, Shao and Chen \cite{r21} also obtained the existence and 
multiplicity of two nontrivial solutions for problem \eqref{lab1.1}.

Motivated by the works mentioned above, especially by  \cite{r16, r19, r21, r33}, 
the purpose of this paper is to study the fractional $p$-Laplacian Kirchhoff type 
problem \eqref{lab1.1} involving steep potential well and more general concave-convex 
nonlinearities. More precisely, we follow the variational methods and use an 
interesting truncation technique as e.g.\ in  \cite{r16} to prove the existence 
and asymptotic behavior of the nontrivial solutions for problem \eqref{lab1.1}. 
Furthermore, we also obtain the nonexistence result of nontrivial solutions to this
problem. It is significantly different from the previous works. When dealing with 
problem \eqref{lab1.1}, we may encounter several difficulties as follows:
\begin{itemize}
\item[(1)] We cannot directly generalize minimum and sufficient for research problem 
\eqref{lab1.1} with
fractional $p$-Laplacian Kirchhoff type due to the set $\Omega\times\Omega$ is strictly 
contained in $\mathbb{R}^{2N}\setminus(\Omega^c\times\Omega^c)$, so we need to 
introduction a new fractional space as \cite{r9, r22, r31} to solve this difficulty;

\item[(2)] Note that $V, ~f$ and $g$ are not required to be radial and 
$\frac{f(x,t)}{|t|^{p-1}}$ may not be increasing in $(-\infty, 0)$
and $(0, +\infty)$ hinders us from borrowing Nehari manifold and fibering methods
such as following \cite{r9, r24, r27}, we will use the truncation technique to
 overcome it;

\item[(3)] Another difficulty is the presence of nonlocal term
$(\int\int_{\mathbb{R}^{2N}}\frac{|u(x)-u(y)|^p}{|x-y|^{N+sp}}\,dx\,dy)(-\Delta)_p^s u$
and the lack of compactness, which means that using the variational method directly 
in a standard way is not feasible. Then we can restore compactness by constructing 
some effective inequalities for $b\to 0$ and $\lambda\to \infty$.
\end{itemize}

To state our main results, we assume that $V$ satisfies:
\begin{itemize}
\item[(A1)] $V(x)\in C(\mathbb{R}^N, \mathbb{R})$ and $V(x)\geq0$ on $\mathbb{R}^N$;

\item[(A2)] there exists ~$V_0>0$ such that the set
$\{V<V_0\}:=\{x\in\mathbb{R}^N: ~V(x)<V_0\}$ is nonempty and has finite Lebesgue measure;

\item[(A3)] $\Omega=$ int$V^{-1}(0)$ is a nonempty open set and has smooth boundary
 with $\overline{\Omega}=V^{-1}(0)$.
\end{itemize}
Considering that $F(x,t)=\int_0^tf(x,\tau)d\tau $ and $G(x,t)=\int_0^tg(x,\tau)d\tau$ 
are the primitives of $f$ and $g$ respectively, we assume the following conditions:
\begin{itemize}
\item[(A4)] $f\in C(\mathbb R^N\times\mathbb R, \mathbb R)$ and $f(x,t)=o(|t|^{p-2}t)$
as $|t|\to  0$ uniformly for $x\in \mathbb R^N$;

\item[(A5)] there exist constants $C>0$ and $q\in(p, p^*_s)$ such that
$|f(x,t)|\leq C(1+|t|^{q-1})$, where $p^*_s=\frac{Np}{N-sp}$ is the fractional 
critical exponent;

\item[(A6)] there exists $\mu>0$ such that $f(x,t)t\geq\mu F(x,t)>0$ for all 
$p<\mu<p^*_s$ and $t\in \mathbb R\setminus\{0\}$;

\item[(A7)] $g\in C(\mathbb R^N\times\mathbb R, \mathbb R)$ and there exists
$h(x)\in L^{\frac{p}{p-r}}(\mathbb R^N, \mathbb R^+)$ such that 
$|g(x,t)|\leq h(x)|t|^{r-1}$ for all $r\in(1,p)$;

\item[(A8)] there exists $\nu\in(1,p)$ such that $\nu G(x,t)\geq tg(x,t)>0$ for all
$t\in \mathbb R \setminus\{0\}$.
\end{itemize}

Now, on the existence and asymptotic behavior of nontrivial solutions for
problem \eqref{lab1.1} we provide the following theorems.


\begin{theorem} \label{thm1.1}
Suppose that {(A1)--(A8)} hold. Then for each fixed $a>0$, there exist
$\widehat{b}, \Lambda, \Pi>0$ such that $b\in(0,\widehat{b})$, $\lambda>\Lambda$ and 
$0<|h|_\frac{p}{p-r}<\Pi$, problem \eqref{lab1.1} possesses at least two different 
nontrivial solutions $u_{b,\lambda}^{(i)}$ $(i=1,2)$.
\end{theorem}

\begin{theorem} \label{thm1.2}
Let $u^{(i)}_{b,\lambda}$ $(i=1,2)$ be two solutions of problem \eqref{lab1.1} 
given by Theorem \ref{thm1.1}. Then for each $b\in(0, \widehat{b})$ and 
$u^{(i)}_{b,\lambda_n}\to  u_b^{(i)}$ in $W^{s,p}(\mathbb{R}^N)$ as 
$\lambda_n\to \infty$ up to a subsequence, where $u_b^{(1)}\neq u_b^{(2)}$ are 
two nontrivial solutions of
\begin{equation}\label{lab1.4}
\begin{gathered}
\Big(a+b\int\int_{\mathbb{R}^{2N}}\frac{|u(x)-u(y)|^p}{|x-y|^{N+sp}}\,dx\,dy\Big)
(-\Delta)_p^su=f(x,u)+g(x,u)\quad \text{in }\Omega,\\
u=0\quad \text{in }\mathbb{R}^N\setminus\Omega.
\end{gathered}
\end{equation}
\end{theorem}

\begin{theorem} \label{thm1.3}
Let $a=1$ and $u^{(i)}_{b,\lambda} (i=1,2)$ be two solutions of problem 
\eqref{lab1.1} given by Theorem \ref{thm1.1}. Then for each fixed $\lambda>\Lambda$
and $u^{(i)}_{b_n,\lambda}\to  u_\lambda^{(i)}$ in $W^{s,p}(\mathbb{R}^N)$ as 
$b_n\to 0$  up to a subsequence, where $u_\lambda^{(1)}\neq u_\lambda^{(2)}$ are 
two nontrivial solutions of
\begin{equation}\label{lab1.5}
\begin{gathered}
(-\Delta)_p^su+\lambda V(x)|u|^{p-2}u=f(x,u)+g(x,u)\ \ \text{in }\mathbb{R}^N,\\
u\in W^{s,p}(\mathbb{R}^N).
\end{gathered}
\end{equation}
\end{theorem}

\begin{theorem} \label{thm1.4}
Let $a=1$ and $u^{(i)}_{b,\lambda} (i=1,2)$ be two solutions of problem 
\eqref{lab1.1} given by Theorem \ref{thm1.1}. Then for 
$u^{(i)}_\lambda\to  \widetilde{u}^{(i)}$ in $W^{s,p}(\mathbb{R}^N)$ as 
$\lambda\to \infty$ and $b\to 0$ up to a subsequence, where 
$\widetilde{u}^{(1)}\neq \widetilde{u}^{(2)}$ are two nontrivial solutions of
\begin{gather*}
(-\Delta)_p^su=f(x,u)+g(x,u)\quad \text{in }\Omega,\\
u=0\quad \text{in }\mathbb{R}^N\setminus\Omega.
\end{gather*}
\end{theorem}

Next, we investigate the nonexistence of nontrivial solutions.

\begin{theorem} \label{thm1.5}
Suppose that {\rm (A1)--(A8)}  hold.
If $a>a_1:=\mathfrak{C}|\overline{\Omega}|^{\frac{p^*_s-p}{p^*_s}}$ and for each 
fixed $\lambda>0$, then there exists a constant $a_2>0$ such that for all $b>0$ and 
$a>\max\{a_1, a_2\}$, problem \eqref{lab1.1} has no nontrivial solution.
\end{theorem}

\begin{remark}  \label{rmk1.6} \rm
Analysis of conditions and regularity of solutions:

(i) $(A1)$--$(A3)$ were first proposed by Bartsch and Wang in \cite{r2},
where $\lambda V(x)$ represents a steep potential well with a bottom of $V^{-1}(0)$, 
and its depth is controlled by $\lambda$.  (A4) and (A5) imply that for any
$\varepsilon> 0$, there exists $C_\varepsilon > 0$ such that
\begin{equation}\label{lab1.6}
|f(x,t)|\leq \varepsilon|t|^{p-1}+C_\varepsilon|t|^{q-1}.
\end{equation}

(ii) Combined with the regularity theory for the fractional $p$-Laplacian equations 
established by Lannizzotto et al.\ \cite{r11, r12}, the kernel function of the operator 
$(-\Delta)_p^s$ satisfies standard singularity and meets the structural requirements 
of the regularity theory. Under assumptions (A1)--(A3), the potential
$V$ ensures the local boundedness of the coefficients in the equation. 
Furthermore, (A4)--(A8) guarantee that $f$ and $g$ satisfy the requirements of
local boundedness and controlled growth for nonlinear terms in the theory. 
Consequently, all nontrivial solutions (denoted by $u$) obtained in this paper are 
globally H\"{o}lder continuous for $i=1,2$:
\begin{itemize}
\item the solutions $u_{b,\lambda}^{(i)}$ in Theorem \ref{thm1.1} and 
 $u_{\lambda}^{(i)}$ in Theorem \ref{thm1.3} satisfy 
 $u\in C^{\alpha}_{\rm loc}(\mathbb{R}^N)$ with some $\alpha\in(0, s]$;
 
\item the solutions $u_{b}^{(i)}$ in Theorem \ref{thm1.2} and $\widetilde{u}^{(i)}$ 
 in Theorem \ref{thm1.4} satisfy $u\in C^{\alpha}(\overline{\Omega})$ with some 
 $\alpha\in(0, s]$, and $\widetilde{u}^{(i)}$ attain the optimal exponent $\alpha=s$;

\item trivial solutions (the nonexistence case in Theorem \ref{thm1.5}) naturally 
enjoy $C^\infty(\mathbb{R}^N)$ regularity.
\end{itemize}

(iii) In the present paper, our main results complement and develop the relevant 
results on fractional $p$-Laplacian Kirchhoff type problems in \cite{r19, r26, r31, r33} 
by assigning more general conditions to the nonlinear terms $f$ and $g$.
\end{remark}

The rest of this article is arranged as follows. 
In the next section, we present some definitions and useful results.
 In Sections $3$ and $4$, we provide proofs for Theorems \ref{thm1.1}--\ref{thm1.4} 
and finally in Section 5 we prove Theorem \ref{thm1.5}.

\section{Preliminaries}

In this section, we will construct a variational framework and provide
preliminary propositions. Before that, we need to some notations. 
$\mathcal{C}, C, C_1, C_2, \cdots$ denote positive constants (possibly different), 
$|\cdot|_\nu$ is the usual norm of the space $L^{\nu}(\mathbb{R}^N)$ for all 
$1\leq \nu\leq \infty$. If the subsequence of sequence $\{u_n\}$ is taken, 
it will be denoted again as $\{u_n\}$. $o(1)$ denotes any quantity which tends 
to zero when $n\to  \infty$. Moreover, define the fractional Sobolev space
$$
W^{s,p}(\mathbb{R}^N)=\{u\in L^p(\mathbb{R}^N): [u]_{s,p}<\infty\}
$$
endowed with the norm
$$
\|u\|_{W^{s,p}(\mathbb{R}^N)}=([u]_{s,p}^p+|u|^p_p)^{1/p}
$$
for
$$
[u]_{s,p}:=\Big(\int\int_{\mathbb{R}^{2N}}\frac{|u(x)-u(y)|^p}{|x-y|^{N+sp}}\,dx\,dy
\Big)^{1/p}.
$$
Now we recall  the embedding relationships of the fractional Sobolev spaces 
into Lebesgue spaces.

\begin{lemma}[see \cite{r5}] \label{L2.1}
Define $\mathcal{D}^{s,p}(\mathbb{R}^N)$ as the closure of
$C_0^\infty(\mathbb{R}^N)$ with respect to $[u]_{s,p}$. Let $s\in(0, 1)$ and 
$p\in[1, \infty)$. Then there exists the optimal Sobolev embedding constant 
$\mathfrak{C}:=\mathfrak{C}(N, s, p)>0$ such that
$$
|u|_{p^*_s}^p\leq\mathfrak{C}[u]_{s,p}^p
$$
for every $u\in\mathcal{D}^{s,p}(\mathbb{R}^N)$. 
Moreover, $W^{s,p}(\mathbb{R}^N)$ is continuously embedded in $L^q(\mathbb{R}^N)$ 
for any $q\in[p, p^*_s]$ and compactly in $L_{\rm loc}^q(\mathbb{R}^N)$ whenever
$q\in[1, p^*_s)$.
\end{lemma}

Set
$$
X=\Big\{u\in W^{s,p}(\mathbb{R}^N):\int_{\mathbb{R}^N}V(x)|u|^pdx<\infty \Big\}
$$
with the norm
$$
\|u\|=\Big(a[u]_{s,p}^p+\int_{\mathbb{R}^N}V(x)|u|^pdx\Big)^{1/p}.
$$
For $\lambda>0$, we also need to provide the following norm
$$
\|u\|_\lambda=\Big(a[u]_{s,p}^p+\int_{\mathbb{R}^N}\lambda V(x)|u|^pdx\Big)^{1/p}.
$$
If $\lambda\geq1$, then $\|u\|\leq\|u\|_\lambda$. Define $X_\lambda=(X, \|u\|_\lambda)$, 
it follows from (A1) and (A2) that
$$
\|u\|^p_{W^{s,p}(\mathbb{R}^N)}
\leq \max \Big\{\frac{\mathfrak{C}|\{V<V_0\}|^{\frac{p^*_s-p}{p^*_s}}+1}{a}, 
\frac{1}{V_0} \Big\}\|u\|^p,
$$
then $X$ is continuously embedded in $W^{s,p}(\mathbb{R}^N)$. 
Let $\Lambda_1=\frac{a}{\mathfrak{C}V_0|\{V<V_0\}|^{\frac{p^*_s-p}{p^*_s}}}$.
As in the idea of  \cite{r23}, for any $\kappa\in[p, p^*_s]$, we also have
\begin{equation}\label{lab2.1}
\begin{gathered}
\int_{\mathbb{R}^N}|u|^\kappa dx
\leq \Big(\frac{\mathfrak{C}}{a}\Big)^{\kappa/p}|\{V<V_0\}
|^{\frac{p^*_s-\kappa}{p^*_s}}
\|u\|^\kappa_\lambda\quad \text{for all}\ \lambda\geq\Lambda_1.
\end{gathered}
\end{equation}
Furthermore, we give additional definitions about fractional space, for more 
details see \cite{r5, r9, r22, r31}. 
Set $\Omega\subset\mathbb{R}^N$ be an open bounded set with smooth boundary, 
$\Omega^c=\mathbb{R}^N\setminus \Omega$, 
$\mathcal{O}=(\Omega^c\times\Omega^c)\subset\mathbb{R}^{2N}$ and 
$\mathcal{Q}=\mathbb{R}^{2N}\setminus \mathcal{O}$. 
It is defined in the usual fractional space $W^{s,p}(\Omega)$ endowed with the norm
$$
\|u\|_{_{W^{s,p}(\Omega)}}
=\Big(\int_{\Omega\times\Omega}\frac{|u(x)-u(y)|^p}{|x-y|^{N+sp}}\,dx\,dy
+\int_{\Omega}|u|^pdx\Big)^{1/p}.
$$
The fractional space $E$ is defined by
$$
E=\Big\{u\in L^p(\Omega): \frac{|u(x)-u(y)|}{|x-y|^{\frac{N}{p}+s}}
\in L^p(\mathcal{Q})\Big\},
$$
and endowed with the norm
$$
\|u\|_E=\Big(\int_\mathcal{Q}\frac{|u(x)-u(y)|^p}{|x-y|^{N+sp}}\,dx\,dy
+\int_\Omega |u|^pdx\Big)^{1/p}.
$$
Moreover, let
$E_0=\{u\in E: u=0\ \text{a.e. in}\ \Omega^c\}$.
Then we have the following results.

\begin{lemma}[see \cite{r9, r31}] \label{L2.2}
The following claims are established:
\begin{itemize}
\item[(i)] if $v\in E$, then $v\in W^{s,p}(\Omega)$ and $\|v\|_{W^{s,p}(\Omega)}
\leq \|v\|_E$;

\item[(ii)] if $v\in E_0$, then $v\in W^{s,p}(\mathbb{R}^N)$ and 
$\|v\|_{W^{s,p}(\Omega)}\leq \|v\|_{W^{s,p}(\mathbb{R}^N)}\leq\|v\|_E$;

\item[(iii)] there exists a constant $\mathfrak{C}:=\mathfrak{C}(N, s, p)>0$ such that 
for any $v\in E_0$,
$$
|v|_{L^{p^*_s}(\Omega)}^p\leq\mathfrak{C}[v]_{s,p}^p;
$$
\item[(iv)] there exists a constant $\widetilde{\mathfrak{C}}
:=\widetilde{\mathfrak{C}}(N, s, p, \Omega)>1$, such that for any $u\in E_0$,
$$
\int_\mathcal{Q}\frac{|u(x)-u(y)|^p}{|x-y|^{N+sp}}\,dx\,dy
\leq\|u\|_E^p\leq\widetilde{\mathfrak{C}}\int_\mathcal{Q}
\frac{|u(x)-u(y)|^p}{|x-y|^{N+sp}}\,dx\,dy,
$$
i.e.
$$
\|u\|_{E_0}=\Big(\int_\mathcal{Q}\frac{|u(x)-u(y)|^p}{|x-y|^{N+sp}}\,dx\,dy
\Big)^{1/p}$$
is a norm on $E_0$ and is equivalent to the norm on $E$;

\item[(v)] $(E_0, \|\cdot\|_{E_0})$ is a reflexive Banach space;

\item[(vi)] let $\{u_n\}\subset E_0$ be a bounded sequence. 
Then, there exists $u\in L^q(\mathbb{R}^N)$ such that up to a
subsequence, $u_n\to  u$ in $L^q(\mathbb{R}^N)$ as $n\to \infty$ for every 
$q\in[1, p^*_s)$.
\end{itemize}
\end{lemma}

Proceeding as in \cite[Lemma 1]{r30} and Lemma $2.2$ in  \cite{r34}, it is easy to 
verify the validity of the following lemma.

\begin{lemma}\label{L2.3} 
Suppose that {\rm (A1)--(A5), (A7)} are satisfied. Then
\begin{itemize}
\item[(i)] $\langle\mathcal{F}'(u), v\rangle=\int_{\mathbb{R}^N}f(x,u)v dx$ 
and $\mathcal{F}': X_\lambda \to  X'_\lambda$ is weakly continuous;

\item[(ii)] $\langle\mathcal{G}'(u), v\rangle=\int_{\mathbb{R}^N}g(x,u)v dx$ and 
$\mathcal{G}': X_\lambda \to  X'_\lambda$ is weakly continuous, where 
$\mathcal{F}(u):=\int_{\mathbb{R}^N}f(x,u)v dx$ and 
$\mathcal{G}(u):=\int_{\mathbb{R}^N}g(x,u)v dx$. 
In particular, if $u_n\rightharpoonup u$ in $X_\lambda$, we have
$$
\lim_{n\to \infty}\int_{\mathbb{R}^N} |g(x,u_n)-g(x,u)|^{\frac{p}{p-1}}dx=0.
$$
\end{itemize}
\end{lemma}

As a consequence, related to problem \eqref{lab1.1} we have the energy functional 
$\mathcal{J}_{b,\lambda}:X_\lambda\to \mathbb{R}$ given by
\[
 \mathcal{J}_{b,\lambda}(u)= \frac{1}{p}\|u\|^p_\lambda+\frac{b}{2p}[u]^{2p}_{s,p}
 -\int_{\mathbb{R}^N}F(x,u)dx-\int_{\mathbb{R}^N}G(x,u)dx.
\]
For all $u\in X_\lambda$ and $v\in C_0^\infty(\mathbb{R}^N)$, it is easy to see that 
$\mathcal{J}_{b,\lambda}\in C^1({X_\lambda},\mathbb{R})$ and
\begin{align*}
\langle \mathcal{J}'_{b,\lambda}(u),v\rangle
&= (a+b[u]^p_{s,p})
\int\int_{\mathbb{R}^{2N}}\frac{|u(x)-u(y)|^{p-2}[u(x)-u(y)][v(x)-v(y)]}
{|x-y|^{N+sp}}\,dx\,dy    \\
&\quad +\int_{\mathbb{R}^N}\lambda V(x)|u|^{p-2}uv\,dx
-\int_{\mathbb{R}^N}f(x,u)v\,dx-\int_{\mathbb{R}^N}g(x, u)v\,dx.
\end{align*}
Clearly, if $u\in X_\lambda$ is a critical point of $\mathcal{J}_{b,\lambda}$, 
then $u$ is a solution of problem \eqref{lab1.1}.

Next, we introduce a variant of the Mountain Pass Theorem  \cite{r7} which is 
considered the Cerami condition. Let $\mathcal{B}$ be a Banach space and 
$\mathcal{J} \in C^1(\mathcal{B}, \mathbb{R})$. 
If a sequence $\{u_n\}\subset \mathcal{B}$ is said to be a Cerami sequence 
(in short $(Ce)_c$ sequence) at the level $c\in \mathbb{R}$ when 
$\mathcal{J}(u_n)\to  c$ and 
$(1+\|u_n\|_\mathcal{B})\|\mathcal{J}'(u_n)\|_{\mathcal{B}^*}\to  0$, where 
$\mathcal{B}^*$ denotes the dual space of $\mathcal{B}$.

\begin{theorem}[see \cite{r7}] \label{thm2.4}
Let $\mathcal{B}$ be a real Banach space. Suppose that 
$\mathcal{J} \in C^1(\mathcal{B}, \mathbb{R})$ satisfies
$$
\max \{\mathcal{J}(0), \mathcal{J}(e)\}
\leq \mu<\eta\leq \inf_{\|u\|_\mathcal{B}=\rho}\mathcal{J}(u)
$$
for some $\eta>\mu, \rho>0$ and $e\in \mathcal{B}$ with $\|e\|_\mathcal{B}>\rho$. 
Let $c\geq \eta$ be characterized by
$$
c=\inf_{\gamma\in\Gamma}\max_{t\in[0,1]}\mathcal{J}(\gamma(t)),
$$
where $\Gamma=\{\gamma\in C([0,1],\mathcal{B}):\gamma(0)=0,\gamma(1)=e\}$ is the set 
of continuous paths joining $0$ and $e$. Then there exists a sequence 
$\{u_n\} \subset \mathcal{B}$ such that
$$
\mathcal{J}(u_n)\to  c\geq \eta\text{ and }
 (1+\|u_n\|_\mathcal{B})\|\mathcal{J}'(u_n)\|_{\mathcal{B}^*}\to  0.
 $$
\end{theorem}

\section{Existence of nontrivial solutions}

In this section, we study the existence of two nontrivial solutions for problem 
\eqref{lab1.1} and provide a proof of Theorem \ref{thm1.1}. For this, we borrow the 
truncation technique. As in \cite{r16}, define the cut-off function 
$\varphi \in C^1([0,\infty),  \mathbb{R})$  satisfying $\varphi(t)=1$ if $t\in[0,1]$, 
$0\leq\varphi(t)\leq1$ if $t\in(1,2)$, $\varphi(t)=0$ if $t\in[2,\infty)$, 
$\max_{t>0}|\varphi'(t)|\leq2$ and $\varphi'(t)\leq0$ for any $t\in(0,\infty)$.
 Furthermore, set
\[
\widehat{\Theta}=\frac{(\mu+r)|h|_{\frac{p}{p-r}}}{\mu r (\Lambda_1 V_0)^{r/p}}
\quad\text{for each }
\beta>\Big(\frac{r\widehat{\Theta}}{2^{\frac{p-r}{p}} p\Theta }\Big)^{\frac{1}{p-r}}
\]
with $\Theta=\frac{\mu-p}{\mu p}$. We move to consider the truncated functional 
$\mathcal{J}^\beta_{b,\lambda}: X_\lambda\to  \mathbb{R}$ defined by
$$
 \mathcal{J}^\beta_{b,\lambda}(u)=\frac{1}{p}\|u\|^p_\lambda+\frac{b}{2p}\varphi\Big( \frac{\|u\|^p_\lambda}{\beta^p}
\Big)[u]^{2p}_{s,p}-\int_{\mathbb{R}^N}F(x,u)dx-\int_{\mathbb{R}^N}G(x,u)dx,
$$
where $\varphi$ is a smooth cut-off function such that
\[
\varphi\Big( \frac{\|u\|^p_\lambda}{\beta^p} \Big)=
\begin{cases}
1, & \|u\|_\lambda \leq \beta, \\
0,& \|u\|_\lambda \geq 2^{1/p}\beta.
\end{cases}
\]
It is easy to see that $\mathcal{J}^\beta_{b,\lambda}$ is of class $C^1$. 
In addition, for all $u, v\in X_\lambda$ we find that
\begin{align*}
\langle (\mathcal{J}^\beta_{b,\lambda})'(u),v\rangle
&= \Big(a\int\int_{\mathbb{R}^{2N}}\frac{|u(x)-u(y)|^{p-2}[u(x)-u(y)]
 [v(x)-v(y)]}{|x-y|^{N+sp}}\,dx\,dy  \\
&\quad +\int_{\mathbb{R}^N}\lambda V(x)|u|^{p-2}uv\,dx\Big) 
\Big[1+\frac{b}{2\beta^p}\varphi'\Big( \frac{\|u\|^p_\lambda}{\beta^p} 
 \Big)[u]_{s,p}^{2p}\Big]  \\
&\quad +b \varphi\Big( \frac{\|u\|^p_\lambda}{\beta^p} \Big)
 [u]_{s,p}^p\int\int_{\mathbb{R}^{2N}}\frac{|u(x)-u(y)|^{p-2}[u(x)-u(y)][v(x)
 -v(y)]}{|x-y|^{N+sp}}\,dx\,dy  \\
&\quad -\int_{\mathbb{R}^N}f(x,u)v\,dx-\int_{\mathbb{R}^N}g(x,u)v\,dx.
\end{align*}

Now we check that the truncated functional $\mathcal{J}^\beta_{b,\lambda}$ 
satisfies the mountain pass geometry.

\begin{lemma}\label{L3.1} 
Suppose that {\rm (A1)--(A7)} are satisfied. Then for all $\lambda\geq\Lambda_1$,
\begin{itemize}
\item[(i)] there exist $\Pi_1, \rho, \eta>0$ (independent of $\beta, b$ and 
$\lambda$) such that for all $0<|h|_{\frac{p}{p-r}}<\Pi_1$,
$$
\inf \{\mathcal{J}^\beta_{b,\lambda}(u): u\in X_\lambda \quad \text{with}\quad
 \|u\|_\lambda=\rho \}\geq\eta;
$$

\item[(ii)] there exist $\widehat{b}>0$ (independent of $\beta$ and $\lambda$) and 
$e\in C_0^\infty(\Omega)$ with $\|e\|_\lambda>\rho$ such that 
$\mathcal{J}^\beta_{b,\lambda}(e)<0$ for all $b\in(0,\widehat{b})$.
\end{itemize}
\end{lemma}

\begin{proof}
Set $\varepsilon=\frac{\Lambda_1 V_0}{2}$. We can use \eqref{lab1.6}, \eqref{lab2.1}, 
(A7) and H\"older inequality to obtain
\begin{align*}
 \mathcal{J}^\beta_{b,\lambda}(u) 
&\geq \frac{1}{p} \|u\|^p_\lambda-\int_{\mathbb{R}^N}F(x,u)dx
 -\int_{\mathbb{R}^N}G(x,u)dx \\
&\geq \frac{1}{p} \|u\|^p_\lambda-\frac{\varepsilon}{p}\int_{\mathbb{R}^N}|u|^pdx
 -\frac{C_\varepsilon}{q} \int_{\mathbb{R}^N}|u|^qdx
 -\frac{1}{r}\int_{\mathbb{R}^N}h(x)|u|^rdx   \\
&\geq  \frac{1}{p}\Big(1-\frac{\varepsilon\mathfrak{C}}{a} |
 \{V<V_0\}|^{\frac{p^*_s-p}{p^*_s}}\Big)\|u\|^p_\lambda
 -\frac{C_\varepsilon}{q}\Big(\frac{\mathfrak{C}}{a}
 \Big)^{q/p}|\{V<V_0\}|^{\frac{p^*_s-q}{p^*_s}}\|u\|^q_\lambda   \\
&\quad -\frac{|h|_{\frac{p}{p-r}}}{r}\Big(\frac{\mathfrak{C}}{a}
 \Big)^{r/p}|\{V<V_0\}|^{\frac{r(p^*_s-p)}{pp^*_s}}\|u\|^r_\lambda  \\
&= \frac{1}{2p}\|u\|^p_\lambda-\frac{C_{\frac{\Lambda_1 V_0}{2}}}{q}
 \Big(\frac{\mathfrak{C}}{a}\Big)^{q/p}
 \Big(\frac{a}{\mathfrak{C}\Lambda_1V_0}\Big)^{\frac{p^*_s-q}{p^*_s-p}}
 \|u\|^q_\lambda -\frac{|h|_{\frac{p}{p-r}}}{r}
 \Big(\frac{1}{\Lambda_1V_0}\Big)^{r/p}\|u\|^r_\lambda    \\
&:= \|u\|^r_\lambda\Big(\frac{1}{2p}\|u\|^{p-r}_\lambda
 -A\|u\|^{q-r}_\lambda-B|h|_{\frac{p}{p-r}}\Big),
\end{align*}
where
$$
A=\frac{C_{\frac{\Lambda_1 V_0}{2}}}{q}\Big(\frac{\mathfrak{C}}{a}\Big)^{q/p}
\Big(\frac{a}{\mathfrak{C}\Lambda_1V_0}\Big)^{\frac{p^*_s-q}{p^*_s-p}}\quad
\text{and}\quad B=\frac{1}{r}\Big(\frac{1}{\Lambda_1V_0}\Big)^{r/p}.
$$
By simple calculation, for each $t\geq0$, we can obtain that
$\Psi(t):=\frac{1}{2p}t^{p-r}-At^{q-r}$
admits a unique maximum point at
$t_0=\Big(\frac{p-r}{2Ap(q-r)}\Big)^{\frac{1}{q-p}}$
and its maximum is
$$
\max_{t\geq0}\Psi(t)=\Psi(t_0)=\frac{q-p}{2p(q-r)}\Big(\frac{p-r}{2Ap(q-r)}
\Big)^{\frac{p-r}{q-p}}>0.
$$
Therefore, it is easy to deduce that there is $t_0:=\rho=\|u\|_\lambda>0$ such that
$$
\mathcal{J}^\beta_{b,\lambda}(u)\geq\eta:=t_0^r\Big(\Psi(t_0)-B|h|_{\frac{p}{p-r}}\Big)>0,
$$
provided that
$$
|h|_{\frac{p}{p-r}}
<\Pi_1:=r(q-p)\Big(\frac{a}{2\mathfrak{C}p(q-r)}\Big)^{\frac{q-r}{q-p}}
\Big(\frac{q(p-r)}{C_{\frac{\Lambda_1 V_0}{2}}}\Big)^{\frac{p-r}{q-p}}
\Big(\frac{\mathfrak{C}\Lambda_1V_0}{a}\Big)^{\frac{(p^*_s-q)(p-r)}{(p^*_s-p)
(q-p)}+\frac{r}{p}},
$$
which completes the proof of (i).

To prove (ii), we define the functional $\mathcal{I}_\lambda: X_\lambda\to  \mathbb{R}$ 
by
$$
\mathcal{I}_\lambda(u)=\frac{1}{p} \|u\|^p_\lambda
-\int_{\mathbb{R}^N}F(x,u)dx-\int_{\mathbb{R}^N}G(x,u)dx.
$$
By (A4)--(A6) and (A8), there exist positive constants $C_i(i=1, 2, \cdots,4)$
such that
\begin{equation}\label{lab3.1}
F(x,t)\geq C_1|t|^\mu- C_2|t|^p \quad \text{and}\quad G(x,t)\geq C_3|t|^\nu-C_4
\end{equation}
for all $(x,t)\in \mathbb R^N\times \mathbb R$.
Let $\phi\in C^\infty_0(\Omega)$ be a fixed smooth positive function. 
Using \eqref{lab3.1}, then one sees that
\begin{align*}
\mathcal{I}_\lambda(l\phi)
&= \frac{al^p}{p}\int_{\mathcal{Q}}\frac{|\phi(x)-\phi(y)|^p}{|x-y|^{N+sp}}\,dx\,dy
 -\int_{\Omega}F(x, l\phi)dx-\int_{\Omega}G(x,l\phi)\,dx  \\
&\leq \frac{al^p}{p}\int_{\mathcal{Q}}\frac{|\phi(x)-\phi(y)|^p}{|x-y|^{N+sp}}\,dx\,dy
 -C_1l^\mu\int_{\Omega}|\phi|^\mu dx+C_2l^p\int_{\Omega}|\phi|^pdx  
 -C_3l^\nu\int_{\Omega}|\phi|^\nu dx+C_4|\Omega|     \\
&\to -\infty\quad \text{as } l\to  +\infty.
\end{align*}
Then, there exists $e\in C^\infty_0(\Omega)$ with $\|e\|_\lambda>\rho$ such 
that $\mathcal{I}_\lambda(e)\leq -1$. 
Then there exists $\widehat{b}>0$ such that
\begin{align*}
 \mathcal{J}^\beta_{b,\lambda}(e)
&=\mathcal{I}_\lambda(e)+\frac{b}{2p}\varphi\Big( \frac{\|e\|^p_\lambda}{\beta^p}
 \Big)\Big(\int_{\mathcal{Q}}\frac{|e(x)-e(y)|^p}{|x-y|^{N+sp}}\,dx\,dy\Big)^2    \\
&\leq -1+\frac{b}{2p}\Big(\int_{\mathcal{Q}}\frac{|e(x)-e(y)|^p}{|x-y|^{N+sp}}\,dx\,dy
\Big)^2   <0
\end{align*}
for all $b\in(0, \widehat{b})$. This completes the proof.
\end{proof}

Now we define the mountain pass value
$$
c_{b,\lambda}^\beta=\inf_{\gamma\in\Gamma} 
\max_{t\in[0,1]}\mathcal{J}^\beta_{b,\lambda}(\gamma(t)),
$$
where $\Gamma=\{\gamma\in C([0,1], X_\lambda):\gamma(0)=0,\gamma(1)=e\}$. 
Then by Lemma \ref{L3.1} and Theorem \ref{thm2.4}, we deduce that for all 
$\lambda\geq\Lambda_1$ and $b\in(0, \widehat{b})$, there exists a 
$(Ce)_{c_{b,\lambda}^\beta}$ sequence $\{u_n\}\subset X_\lambda$ such that
\begin{equation}\label{lab3.2}
 \mathcal{J}^\beta_{b,\lambda}(u_n)\to  c_{b,\lambda}^\beta\geq \eta>0\quad
  \text{and}\quad (1+\|u_n\|_\lambda)\|(\mathcal{J}^\beta_{b,\lambda})'(u_n)
  \|_{X_\lambda^*}\to  0.
\end{equation}

In the following lemma, we provide an estimate on the upper bound of 
$c_{b,\lambda}^\beta$ which plays a key role in the use of truncation technique.

\begin{lemma}\label{L3.2} 
Suppose that {\rm (A1)--(A6), (A8)} are satisfied. Then for all
$\lambda\geq\Lambda_1$ and $b\in(0, \widehat{b})$, there exists $M>0$ 
(independent of $\beta, b$ and $\lambda$) such that $c_{b,\lambda}^\beta\leq M$.
\end{lemma}

\begin{proof}
For each $0<e\in C^\infty_0(\Omega)$, we have
\begin{align*}
 \mathcal{J}^\beta_{b,\lambda}(te)
&= \frac{t^p}{p} \|e\|^p_\lambda+\frac{bt^{2p}}{2p}\varphi
 \Big( \frac{t^p\|e\|^p_\lambda}{\beta^p} \Big)[e]^{2p}_{s,p}-
 \int_{\mathbb{R}^N} F(x,te)dx-\int_{\mathbb{R}^N} G(x,te)dx  \\
&\leq  \frac{at^p}{p}\int_{\mathcal{Q}}\frac{|e(x)-e(y)|^p}{|x-y|^{N+sp}}\,dx\,dy
 +\frac{\widehat{b}t^{2p}}{2p}
 \Big(\int_{\mathcal{Q}}\frac{|e(x)-e(y)|^p}{|x-y|^{N+sp}}\,dx\,dy\Big)^2  \\
&\quad -C_1t^\mu\int_{\Omega}|e|^\mu dx+C_2t^p\int_{\Omega}|e|^pdx
 -C_3t^\nu\int_{\Omega}|e|^\nu dx+C_4|\Omega|.
\end{align*}
Consequently, there exists a constant $M>0$ such that
$$
c_{b,\lambda}^\beta \leq \max_{t\in[0,1]}\mathcal{J}^\beta_{b,\lambda}(te)\leq M.
$$
This completes the proof.
\end{proof}

Next, up to a subsequence, we show that the sequence $\{u_n\}$ given by \eqref{lab3.2} 
satisfies $\|u_n\|_{\lambda}\leq\beta$ for the selected appropriate $\beta>0$ which is 
defined at the beginning of this section.

\begin{lemma}\label{L3.3}
Suppose that {\rm (A1)--(A3), (A6)--(A8)} are satisfied, and let
$\beta=\big(\frac{M+1}{\Theta}\big)^{1/p}$.
If $\{u_n\}\subset X_\lambda$ is a sequence satisfying \eqref{lab3.2}, up to a 
subsequence, then there exists $\Pi_2>0$ such that for all $b\in(0, \widehat{b})$ 
and $0<|h|_{\frac{p}{p-r}}<\Pi_2$, we have $\|u_n\|_\lambda\leq \beta$ for all 
$\lambda\geq\Lambda_1$.
\end{lemma}

\begin{proof}
We first prove that $\|u_n\|_\lambda\leq 2^{1/p}\beta$ as $n\to \infty$.
Suppose by contradiction that, there exists a subsequence of $\{u_n\}$, still 
denoted by $\{u_n\}$, such that $\|u_n\|_\lambda> 2^{1/p}\beta$.
We introduce the function $\Phi: [0,+\infty)\to  \mathbb{R}$ defined by
$\Phi(t)=\Theta t^p-\widehat{\Theta}t^r$.

A direct calculation shows that there exists a unique number
$t_0=\big(\frac{r\widehat{\Theta}}{p\Theta}\big)^\frac{1}{p-r}>0$
such that $\min_{t\geq0}\Phi(t)=\Phi(t_0)$. 
Thus, using (A6)--(A8), \eqref{lab2.1} and H\"older inequality lead to
\begin{align*}
 c_{b,\lambda}^\beta
&= \lim_{n\to \infty} \Big[\mathcal{J}^\beta_{b,\lambda}(u_n)-\frac{1}{\mu} 
 \langle (\mathcal{J}^\beta_{b,\lambda})'(u_n), u_n\rangle \Big] \\
&= \lim_{n\to \infty} \Big[\Big(\frac{1}{p}-\frac{1}{\mu}-\frac{b}{2\mu\beta^p}
 \varphi'\Big( \frac{\|u_n\|^p_\lambda}{\beta^p}\Big) [u_n]^{2p}_{s,p}\Big)
 \|u_n\|^p_\lambda+\Big(\frac{b}{2p}-\frac{b}{\mu}\Big)\varphi
 \Big( \frac{\|u_n\|^p_\lambda}{\beta^p}\Big) [u_n]^{2p}_{s,p}  \\
&\quad -\int_{\mathbb{R}^N}\Big(F(x,u_n)-\frac{1}{\mu}f(x,u_n)u_n\Big)dx
 -\int_{\mathbb{R}^N}\Big(G(x,u_n)-\frac{1}{\mu}g(x,u_n)u_n\Big)                        dx \Big]   \\
&\geq \lim_{n\to \infty} \Big[\Big(\frac{\mu-p}{\mu p}
 -\frac{b}{2\mu\beta^p} \varphi'\Big( \frac{\|u_n\|^p_\lambda}{\beta^p}\Big)
  [u_n]^{2p}_{s,p}\Big)\|u_n\|^p_\lambda-\frac{b(2p-\mu)}{2\mu p}
  \varphi\Big( \frac{\|u_n\|^p_\lambda}{\beta^p}\Big) [u_n]^{2p}_{s,p}  \\
&\quad -\int_{\mathbb{R}^N} \Big(|G(x,u_n)|+\frac{1}{\mu}|g(x,u_n)u_n|\Big)dx \Big]  \\
&\geq \lim_{n\to \infty} \Big[ \Big(\frac{\mu-p}{\mu p} -\frac{b}{2\mu\beta^p} 
 \varphi'\Big( \frac{\|u_n\|^p_\lambda}{\beta^p}\Big) [u_n]^{2p}_{s,p}\Big)
 \|u_n\|^p_\lambda-\frac{b(2p-\mu)}{2\mu p}\varphi
 \Big( \frac{\|u_n\|^p_\lambda}{\beta^p}\Big) [u_n]^{2p}_{s,p}  \\
&\quad -\frac{(\mu+r)|h|_{\frac{p}{p-r}}}{\mu r}\Big(\frac{1}{\Lambda_1V_0}
 \Big)^{r/p}\|u_n\|^r_\lambda  \Big]  \\
&= \lim_{n\to \infty} \Big[\frac{\mu-p}{\mu p}\|u_n\|^p_\lambda
 -\frac{(\mu+r)|h|_{\frac{p}{p-r}}}{\mu r (\Lambda_1 V_0)^{r/p}}
 \|u_n\|^r_\lambda-\frac{b}{2\mu\beta^p}
\varphi'\Big(\frac{\|u_n\|^p_\lambda}{\beta^p}\Big) [u_n]^{2p}_{s,p}\|u_n\|^p_\lambda  \\
&\quad -\frac{b(2p-\mu)}{2\mu p}\varphi\Big( \frac{\|u_n\|^p_\lambda}{\beta^p}\Big)
  [u_n]^{2p}_{s,p}\Big]  \\
&\geq 2M+\frac{3}{2}+\Big[\frac{1}{2}-\widehat{\Theta} 
 \Big(\frac{2(M+1)}{\Theta}\Big)^{r/p} \Big]  \\
&> 2M+\frac{3}{2},
\end{align*}
provided that
\begin{align*}
|h|_{\frac{p}{p-r}}
<\Pi_2:=\min\Big\{\frac{\mu r}{\mu+r} \Big[\frac{\Lambda_1 V_0(\mu-p)}{2^{\frac{p+r}{r}}
 \mu p(M+1)}\Big]^{r/p},
\frac{2\mu p(M+1)}{\mu+r}\Big[\frac{\Lambda_1 V_0(\mu-p)}{2\mu p(M+1)}
 \Big]^{r/p}\Big\},
\end{align*}
which is a contradiction by Lemma \ref{L3.2}.

To finish the proof of the lemma, let us argue by contradiction. 
Assuming that there is no subsequence of $\{u_n\}$ which is uniformly bounded by 
$\beta$. Then we infer that $\beta<\|u_n\|_\lambda \leq 2^{1/p}\beta$ as
$n\to  \infty$. Using a calculation method similar to the above and borrowing the 
fact that $\varphi$ is non-increasing, we know that
\begin{align*}
 c_{b,\lambda}^\beta
&= \lim_{n\to \infty} \Big[\mathcal{J}^\beta_{b,\lambda}(u_n)-\frac{1}{\mu} \langle (\mathcal{J}^\beta_{b,\lambda})'(u_n), u_n\rangle \Big]    \\
&\geq \lim_{n\to \infty} \Big[\frac{\mu-p}{\mu p}\|u_n\|^p_\lambda-\frac{(\mu+r)|h|_{\frac{p}{p-r}}}{\mu r (\Lambda_1 V_0)^{r/p}}\|u_n\|^r_\lambda   \\
&-\frac{b}{2\mu\beta^p} \varphi'\Big( \frac{\|u_n\|^p_\lambda}{\beta^p}\Big) [u_n]^{2p}_{s,p}\|u_n\|^p_\lambda
-\frac{b(2p-\mu)}{2\mu p}\varphi\Big( \frac{\|u_n\|^p_\lambda}{\beta^p}\Big) [u_n]^{2p}_{s,p}\Big]  \\
&\geq \liminf_{n\to \infty} \Big[\frac{\mu-p}{\mu p}\|u_n\|^p_\lambda-\frac{(\mu+r)|h|_{\frac{p}{p-r}}}{\mu r (\Lambda_1 V_0)^{r/p}}\|u_n\|^r_\lambda-\frac{b|2p-\mu|}{2a^2\mu p} \|u_n\|^{2p}_\lambda\Big]  \\
&\geq \Theta\beta^p-\widehat{\Theta}\Big(2^{1/p}\beta\Big)^r-\frac{b|2p-\mu|}{2a^2\mu p}\Big(2^{1/p}\beta\Big)^{2p}   \\
&= M+\frac{1}{2}+\Big[\frac{1}{2}-\widehat{\Theta} \Big(\frac{2(M+1)}{\Theta}\Big)^{r/p}\Big]-\frac{2b|2p-\mu|}{\mu p} \Big(\frac{M+1}{a\Theta}\Big)^2,
\end{align*}
provided that $0<|h|_{\frac{p}{p-r}}<\Pi_2$, there exists a constant $\widehat{b}>0$ 
such that for each fixed $a>0$ and $b\in(0, \widehat{b})$, this is a contradiction 
by choosing $\widehat{b}$ as small enough. This completes the proof.
\end{proof}

\begin{remark}  \label{rmk3.4} \rm
If the sequence $\{u_n\}\subset X_\lambda$ is obtained in Lemma \ref{L3.3}, 
then $\{u_n\}$ is also a bounded $(Ce)$ sequence of $\mathcal{J}_{b,\lambda}$ satisfying 
$\|u_n\|_\lambda\leq \beta$. By the definition of the truncation function $\varphi$, 
one can see that
\[
\mathcal{J}_{b,\lambda}(u_n)\to  c_{b,\lambda}^\beta\quad \text{and}\quad
 (1+\|u_n\|_\lambda)\|\mathcal{J}_{b,\lambda}'(u_n)\|_{X_\lambda^*}\to  0.
\]
\end{remark}

\begin{lemma}\label{L3.5}
Suppose that {\rm (A1)--(A3), (A6)--(A8)} are satisfied.
Let $\kappa\in[p, p^*_s)$ and $\{u_n\}$ be a bounded $(Ce)_{c_{b,\lambda}^\beta}$ 
sequence. Then there is a subsequence again denoted as $\{u_n\}$ such that for each 
$\varepsilon>0$, there exists $R_\varepsilon>0$ such that for all 
$R\geq R_\varepsilon$,
$$
\limsup_{n\to \infty}\int_{B_n\setminus B_R}|u_n|^\kappa dx\leq \varepsilon,
$$
where $B_k=\{x\in\mathbb{R}^N: |x|\leq k\}$.
\end{lemma}

The  proof of the above lemma can be found in  \cite{r15}, so we omit it here.

\begin{lemma}\label{L3.6}
Suppose that {\rm (A4)} and {\rm (A5)} are satisfied.
If $u_n \rightharpoonup u$ in $W^{s,p}(\mathbb{R}^N)$ and $v_n:=u_n-u$, then
$$
\lim_{n\to \infty}\int_{\mathbb{R}^N} |[f(x,u_n)-f(x,v_n)-f(x,u)]\psi|dx=0
$$
uniformly in $\psi\in W^{s,p}(\mathbb{R}^N)$ with 
$\|\psi\|_{W^{s,p}(\mathbb{R}^N)}\leq1$.
\end{lemma}

\begin{proof}
The proof is similar to the proof in \cite{r6}. 
For the convenience of readers, here we provide the process of proof. 
Let $\xi: [0, +\infty)\to [0,1]$ be a smooth cut-off function satisfying $\xi(t)=1$ 
if $|t|\leq1$ and $\xi(t)=0$ if $|t|\geq p$. 
We define $\bar{u}_n(x)=\xi(\frac{p|x|}{n})u(x)$. Obviously,
\begin{equation}\label{lab3.3}
\|\bar{u}_n-u\|_\lambda\to 0 \quad \text{as } n\to \infty.
\end{equation}
Going if necessary to a subsequence, define $w_n=u-\bar{u}_n$. 
It follows from \eqref{lab3.3} and any $\psi\in W^{s,p}(\mathbb{R}^N)$ that
\begin{equation}\label{lab3.4}
\lim_{n\to \infty}\sup_{\|\psi\|_{W^{s,p}(\mathbb{R}^N)}\leq1}
\int_{\mathbb{R}^N} [f(x,\bar{u}_n)-f(x,u)]\psi dx=0.
\end{equation}
Moreover, by (A4), (A5), Lemma \ref{L3.5} and \cite[Lemma 3.3]{r15}, a standard
argument shows that
\begin{equation}\label{lab3.5}
\lim_{n\to \infty}\sup_{\|\psi\|_{W^{s,p}(\mathbb{R}^N)}\leq1}
\int_{\mathbb{R}^N} [f(x,u_n)-f(x,u_n-\bar{u}_n)-f(x,\bar{u}_n)]\psi dx=0.
\end{equation}
Next, we prove that
\begin{equation}\label{lab3.6}
\lim_{n\to \infty}\sup_{\|\psi\|_{W^{s,p}(\mathbb{R}^N)}\leq1}
\int_{\mathbb{R}^N} [f(x,v_n+w_n)-f(x,v_n)]\psi dx=0.
\end{equation}
We define $\bar{F}(x,0)=0$ and
$\bar{F}(x,t)=\frac{f(x,t)}{|t|^{p-2}t}$ if $t\neq0.$
By (A4) and (A5), $\bar{F}$ is continuous at $t=0$ and
$|\bar{F}(x,t)|\leq C_1(1+|t|^{q-p})$ for all $(x,t)\in\mathbb{R}^N\times\mathbb{R}$. 
For any $\sigma>0$, we set $A^\sigma_n=\{x\in \mathbb{R}^N: |v_n(x)|\leq \sigma \}$ 
and $B^\sigma_n=\mathbb{R}^N\setminus A^\sigma_n$. Then the Lebesgue measure
$$
|B^\sigma_n|\leq\frac{1}{\sigma^q}\int_{B^\sigma_n}|v_n|^qdx
\leq\frac{C_2}{\sigma^q}\to 0\ \ \text{as}\ \sigma\to \infty.
$$
Hence, for some $\bar{\sigma}>0$ and any $\varepsilon>0$, by (A4), (A5),
H\"older inequality and the boundedness of $\{w_n\}$, we have
\begin{equation}\label{lab3.7}
\Big|\int_{B^{\bar{\sigma}}_n} [f(x,v_n+w_n)-f(x,v_n)]\psi dx\Big|\leq\varepsilon
\end{equation}
uniformly in $\|\psi\|_{W^{s,p}(\mathbb{R}^N)}\leq1$. 
By the uniformly continuity of $\bar{F}$ on 
$\mathbb{R}^N\times[-\bar{\sigma}, \bar{\sigma}]$, there exists 
$\delta:=\delta(\varepsilon, \bar{\sigma})>0$ such that $|w|<\delta$ and
\begin{equation}\label{lab3.8}
|\bar{F}(x,t+w)-\bar{F}(x,t)|\leq\varepsilon\quad \text{for all }
  (x,t)\in\mathbb{R}^N\times[-\bar{\sigma}, \bar{\sigma}].
\end{equation}
Similarly, we set $C^\delta_n=\{x\in\mathbb{R}^N: |w_n(x)|\leq\delta \}$ and 
$D^\delta_n=\mathbb{R}^N\setminus C^\delta_n$. 
Since $|A^{\bar{\sigma}}_n\bigcap D^\delta_n|\to 0$ as $n\to \infty$, 
it is easy to know that there exists $N_1\in \mathbb{N}^+$ such that $n\geq N_1$ and
\begin{equation}\label{lab3.9}
\Big|\int_{A^{\bar{\sigma}}_n\bigcap D^\delta_n} [f(x,v_n+w_n)
-f(x,v_n)]\psi dx\Big|\leq\varepsilon
\end{equation}
uniformly in $\|\psi\|_{W^{s,p}(\mathbb{R}^N)}\leq1$. 
By using \eqref{lab3.3}, there exists $N_2\in \mathbb{N}^+$ ensures that both 
$|w_n|_p<\varepsilon$ and $|w_n|_q<\varepsilon$ hold for any $n\geq N_2\geq N_1$. 
Thus, from \eqref{lab3.8} we infer that
\begin{equation}\label{lab3.10}
\Big|\int_{A^{\bar{\sigma}}_n\bigcap C^\delta_n} [\bar{F}(x,v_n+w_n)
-\bar{F}(x,v_n)]|v_n|^{p-2}v_n\psi dx\Big|\leq C_3\varepsilon
\end{equation}
uniformly in $\|\psi\|_{W^{s,p}(\mathbb{R}^N)}\leq1$. 
Based on \eqref{lab3.8}--\eqref{lab3.10}, the boundedness of $\{v_n\}$ and 
$\{w_n\}$, H\"older inequality and the  inequalities
\begin{equation}\label{lab3.11}
\begin{gathered}
|a+b|^\alpha\leq C_{\alpha}(|a|^\alpha+|b|^\alpha), \quad \alpha\in(0, +\infty),  \\
\|a|^{\alpha-2}a-|b|^{\alpha-2}b|\leq C_\alpha(|a|+|b|)^{\alpha-2}|a-b|, \quad
 \alpha\in(2, +\infty), 
\end{gathered}
\end{equation}
one has
\begin{align*}
&\Big|\int_{A^{\bar{\sigma}}_n} [f(x,v_n+w_n)-f(x,v_n)]\psi dx\Big|   \\
&\leq\Big|\int_{A^{\bar{\sigma}}_n\bigcap C^\delta_n} [f(x,v_n+w_n)-f(x,v_n)]\psi dx\Big|  \\
&\ \ \ \ +\Big|\int_{A^{\bar{\sigma}}_n\bigcap D^\delta_n} [f(x,v_n+w_n)-f(x,v_n)]\psi dx\Big|   \\
&\leq\Big|\int_{A^{\bar{\sigma}}_n\bigcap D^\delta_n} [f(x,v_n+w_n)-f(x,v_n)]\psi dx\Big|   \\
&\ \ \ \ +\Big|\int_{A^{\bar{\sigma}}_n\bigcap C^\delta_n} [\bar{F}(x,v_n+w_n)-\bar{F}(x,v_n)]|v_n|^{p-2}v_n\psi dx\Big|   \\
&\ \ \ \ +\Big|\int_{A^{\bar{\sigma}}_n\bigcap C^\delta_n} \bar{F}(x,v_n+w_n)[|v_n+w_n|^{p-2}(v_n+w_n)-|v_n|^{p-2}v_n]\psi dx\Big|   \\
&\leq\Big|\int_{A^{\bar{\sigma}}_n\bigcap D^\delta_n} [f(x,v_n+w_n)-f(x,v_n)]\psi dx\Big|   \\
&\ \ \ \ +\Big|\int_{A^{\bar{\sigma}}_n\bigcap C^\delta_n} [\bar{F}(x,v_n+w_n)-\bar{F}(x,v_n)]|v_n|^{p-2}v_n\psi dx\Big|   \\
&\ \ \ \ +C_4\int_{\mathbb{R}^N} (1+|v_n+w_n|^{q-p})(|v_n|^{p-2}|w_n|+|w_n|^{p-1})|\psi| dx   \\
&\leq \varepsilon+C_3\varepsilon+C_5(\varepsilon+\varepsilon^{q-1}+\varepsilon^{p-1}+\varepsilon^{q-p+1})
\end{align*}
uniformly in $\|\psi\|_{W^{s,p}(\mathbb{R}^N)}\leq1$. 
This, together with \eqref{lab3.7}, implies that \eqref{lab3.6}.  Hence, it follows 
from \eqref{lab3.4}--\eqref{lab3.6} that Lemma holds.
\end{proof}

We are now preparing to prove the following compactness conditions of 
$\mathcal{J}_{b,\lambda}$.

\begin{lemma}\label{L3.7}
Suppose that {\rm (A1)--(A5), (A7)} are satisfied, and let
$\beta=\big(\frac{M+1}{\Theta}\big)^{1/p}$. Then there exists $\Lambda>0$
such that, for each $b\in(0, \widehat{b})$ and $\lambda>\Lambda$, 
if $\{u_n\}\subset X_\lambda$ is a sequence satisfying \eqref{lab3.2}, 
then $\{u_n\}$ has a convergent subsequence in $X_\lambda$.
\end{lemma}

\begin{proof}
By Lemma \ref{L3.3} and Remark \ref{rmk3.4}, we see that the bounded 
$(Ce)_{c_{b,\lambda}^\beta}$ sequence $\{u_n\}$ satisfying $\|u_n\|_\lambda\leq \beta$. 
Then there exist $u\in X_\lambda$ and $L>0$ such that
\begin{equation}\label{lab3.12}
\begin{gathered}
u_n\rightharpoonup u \quad \text{weakly in } X_\lambda, \\
u_n\to  u  \quad \text{strongly in } L^\kappa_{\rm loc}(\mathbb{R}^N) \quad
\text{for all } \kappa\in[p,p_s^*), \\
u_n\to  u  \quad \text{a.e. in } \mathbb{R}^N, \\
[u]^p_{s,p}\leq \lim_{n\to \infty}[u_n]^p_{s,p}=L^p.
\end{gathered}
\end{equation}
Since $X_\lambda$ is continuously embedded in $W^{s,p}(\mathbb{R}^N)$, 
there exists a constant $C>0$ such that $\|u_n\|_{W^{s,p}(\mathbb{R}^N)}\leq C\beta$. 
Then, $u_n\rightharpoonup u$ in $W^{s,p}(\mathbb{R}^N)$ implies that
\begin{equation}\label{lab3.13}
\begin{split}
&\lim_{n\to \infty}\int\int_{\mathbb{R}^{2N}}
 \frac{|u_n(x)-u_n(y)|^{p-2}[u_n(x)-u_n(y)][v(x)-v(y)]}{|x-y|^{N+sp}}\,dx\,dy\\
&=\int\int_{\mathbb{R}^{2N}}
 \frac{|u(x)-u(y)|^{p-2}[u(x)-u(y)][v(x)-v(y)]}{|x-y|^{N+sp}}\,dx\,dy..
\end{split}
\end{equation}
Combining Lemma \ref{L2.3} and H\"older inequality, for all $v\in X_\lambda$, 
we obtain
\begin{gather}\label{lab3.14}
\lim_{n\to \infty}\int_{\mathbb{R}^N}f(x, u_n)v\,dx=\int_{\mathbb{R}^N}f(x, u)v\,dx,\\
\label{lab3.15}
\lim_{n\to \infty}\int_{\mathbb{R}^N}g(x, u_n)v\,dx=\int_{\mathbb{R}^N}g(x, u)v\,dx.
\end{gather}
Then, by \eqref{lab3.12}--\eqref{lab3.15}, for any $v\in C^\infty_0(\mathbb{R}^N)$, 
we have
\begin{align*}
o(1)&= \langle \mathcal{J}_{b,\lambda}'(u_n), v \rangle  \\
&\to  (a+bL^p)\int\int_{\mathbb{R}^{2N}}\frac{|u(x)-u(y)|^{p-2}[u(x)-u(y)]
 [v(x)-v(y)]}{|x-y|^{N+sp}}\,dx\,dy  \\
&\quad +\int_{\mathbb{R}^N}\lambda V(x)|u|^{p-2}uv\,dx 
 -\int_{\mathbb{R}^N}f(x,u)v\,dx-\int_{\mathbb{R}^N}g(x,u)v\,dx \quad
  \text{as } n\to \infty,
\end{align*}
which yields
\begin{equation}\label{lab3.16}
\|u\|^p_\lambda+bL^p[u]^p_{s,p} -\int_{\mathbb{R}^N}f(x,u)u\,dx
-\int_{\mathbb{R}^N}g(x,u)u\,dx=0.
\end{equation}
Set $v_n=u_n-u$. By $\|u_n\|_\lambda\leq \beta$ one has
\[
\|u\|_\lambda \leq \liminf_{n\to \infty}\|u_n\|_\lambda\leq\beta,
\]
leading to
\begin{equation}\label{lab3.17}
\|v_n\|_\lambda=\|u_n-u\|_\lambda\leq 2\beta.
\end{equation}
From the Br\'{e}zis-Lieb Lemma \cite{r3} we know that
\begin{equation}\label{lab3.18}
\|v_n\|^p_\lambda=\|u_n\|^p_\lambda-\|u\|^p_\lambda+o(1).
\end{equation}
Moreover, in view of (A2), \eqref{lab3.12} and H\"older inequality,
for any $\kappa\in[p, p^*_s)$, one has
\begin{align*}
\int_{\mathbb{R}^N}|v_n|^pdx
&=\int_{\{V\geq V_0\}}|v_n|^pdx+\int_{\{V<V_0\}}|v_n|^pdx      \\
&\leq \frac{1}{\lambda V_0}\int_{\mathbb{R}^N}\lambda V|v_n|^pdx
 +|{\{V<V_0\}}|^{\frac{\kappa-p}{\kappa}}
 \Big(\int_{\{V<V_0\}}|v_n|^\kappa dx\Big)^{\frac{p}{\kappa}}   \\
&\leq \frac{1}{\lambda V_0}\|v_n\|^p_\lambda+o(1).
\end{align*}
This implies
\begin{equation}\label{lab3.19}
\begin{split}
\int_{\mathbb{R}^N}|v_n|^q dx
&=  \int_{\mathbb{R}^N}\Big(|v_n|^{\frac{p(p_s^*-q)}{p_s^*-p}}
 |v_n|^{\frac{p_s^*(q-p)}{p_s^*-p}}\Big)dx \\
&\leq  \Big(\int_{\mathbb{R}^N}|v_n|^pdx\Big)^{\frac{p_s^*-q}{p_s^*-p}} 
 \Big(\int_{\mathbb{R}^N}|v_n|^{p_s^*} dx \Big)^{\frac{q-p}{p_s^*-p}}  \\
&\leq   \Big(\frac{1}{\lambda V_0}\|v_n\|^p_\lambda \Big)^{\frac{p_s^*-q}{p_s^*-p}}
 \Big(\mathfrak{C}\int\int_{\mathbb{R}^{2N}}\frac{|v_n(x)-v_n(y)|^p}{|x-y|^{N+sp}}\,dx\,dy\Big)^{\frac{p_s^*(q-p)}{p(p_s^*-p)}}+o(1)  \\
&\leq   \Big(\frac{1}{\lambda V_0}\Big)^{\frac{p_s^*-q}{p_s^*-p}}
 \Big(\frac{\mathfrak{C}}{a}\Big)^{\frac{p_s^*(q-p)}{p(p_s^*-p)}}\|v_n\|^q_\lambda+o(1).
\end{split}
\end{equation}
By Lemma \ref{L3.6}, $v_n\rightharpoonup0$ in $W^{s,p}(\mathbb{R}^N)$ and $v_n\to  0$ 
in $L^\kappa_{\rm loc}(\mathbb{R}^N)$ for all $\kappa\in[p,p_s^*)$, we obtain
\begin{equation}\label{lab3.20}
\begin{split}
\int_{\mathbb{R}^N} [f(x,u_n)u_n-f(x,u)u]dx 
&=\int_{\mathbb{R}^N} [f(x,u_n)-f(x,v_n)-f(x,u)u_n]dx \\
&\quad +\int_{\mathbb{R}^N} [f(x,v_n)v_n+f(x,v_n)u+f(x,u)v_n]dx  \\
&=\int_{\mathbb{R}^N} f(x,v_n)v_n\,dx+o(1).
\end{split}
\end{equation}
Based on \eqref{lab1.6}, \eqref{lab2.1}, \eqref{lab3.17}, \eqref{lab3.19}, 
H\"older inequality and let $\varepsilon=\frac{\Lambda_1 V_0}{2}$, one has
\begin{equation}\label{lab3.21}
\begin{split}
&\int_{\mathbb{R}^N} f(x,v_n)v_n\,dx \\
&\leq \frac{\Lambda_1 V_0}{2}\int_{\mathbb{R}^N}|v_n|^pdx+C_{\frac{\Lambda_1 V_0}{2}}\int_{\mathbb{R}^N}|v_n|^qdx \\
&\leq \frac{\Lambda_1 V_0}{2}\int_{\mathbb{R}^N}|v_n|^pdx
+C_{\frac{\Lambda_1 V_0}{2}}\Big(\int_{\mathbb{R}^N}|v_n|^qdx
 \Big)^{\frac{q-p}{q}}\Big(\int_{\mathbb{R}^N}|v_n|^qdx\Big)^{\frac{p}{q}}  \\
&\leq \Big[\frac{1}{2}+C_{\frac{\Lambda_1 V_0}{2}}\Big((2\beta)^q
 \Big(\frac{1}{\lambda V_0}\Big)^{\frac{p_s^*-q}{p_s^*-p}}
 \Big(\frac{\mathfrak{C}}{a}\Big)^{\frac{p_s^*(q-p)}{p(p_s^*-p)}}\Big)
^{\frac{q-p}{q}}\frac{\mathfrak{C}}{a} 
 \Big(\frac{a}{\mathfrak{C}\Lambda_1V_0}\Big)^{\frac{p(p^*_s-q)}{q(p^*_s-p)}}
 \Big]\|v_n\|^p_\lambda+o(1).
\end{split}
\end{equation}
Furthermore, by applying Lemma \ref{L2.3}, we  obtain 
\begin{equation}\label{lab3.22}
\begin{split}
&\int_{\mathbb{R}^N}[g(x,u_n)u_n-g(x,u)u]dx  \\
&=\int_{\mathbb{R}^N} g(x,u)(u_n-u)dx+\int_{\mathbb{R}^N} [g(x,u_n)-g(x,u)u_n]dx \\
&\leq\int_{\mathbb{R}^N} g(x,u)v_n\,dx+\Big(\int_{\mathbb{R}^N} |g(x,u_n)-g(x,u)|^{\frac{p}{p-1}}dx\Big)^{\frac{p-1}{p}}|u_n|_p \\
&\to  0\quad \text{as }  n\to \infty.
\end{split}
\end{equation}
Observe that $\langle \mathcal{J}_{b,\lambda}'(u_n), u_n\rangle=o(1)$, 
then from \eqref{lab3.12}, \eqref{lab3.16}, \eqref{lab3.18} and 
\eqref{lab3.20}--\eqref{lab3.22}, we have
\begin{align*}
o(1)&= \|v_n\|_\lambda^p+bL^{2p}-bL^p[u]^p_{s,p}
-\int_{\mathbb{R}^N}[f(x,u_n)u_n-f(x,u)u]dx  \\
&-\int_{\mathbb{R}^N}[g(x,u_n)u_n-g(x,u)u]dx +o(1)    \\
&\geq  \|v_n\|_\lambda^p-\int_{\mathbb{R}^N}[f(x,u_n)u_n-f(x,u)u]dx+o(1)  \\
&\geq  \Big[\frac{1}{2}-C_{\frac{\Lambda_1 V_0}{2}}\Big((2\beta)^q
 \Big(\frac{1}{\lambda V_0}\Big)^{\frac{p_s^*-q}{p_s^*-p}}
 \Big(\frac{\mathfrak{C}}{a}\Big)^{\frac{p_s^*(q-p)}{p(p_s^*-p)}}\Big)^{\frac{q-p}{q}}
 \frac{\mathfrak{C}}{a} \Big(\frac{a}{\mathfrak{C}\Lambda_1V_0}
 \Big)^{\frac{p(p^*_s-q)}{q(p^*_s-p)}}\Big]\|v_n\|_\lambda^p+o(1),
\end{align*}
which yields that there exists
$$
\Lambda_2=\frac{(2\beta)^{\frac{q(p^*_s-p)}{p^*_s-q}}}{V_0}
 |\{V<V_0\}|^{\frac{p(p^*_s-p)}{p^*_s(q-p)}}
\Big(2C_{\frac{\Lambda_1 V_0}{2}}\Big)^{\frac{q(p^*_s-p)}{(q-p)(p^*_s-q)}}
\Big(\frac{\mathfrak{C}}{a}\Big)^{\frac{pq(p_s^*-p)+p_s^*(q-p)^2}{p(q-p)(p_s^*-q)}}
$$
such that $v_n\to 0$ strongly in $X_\lambda$ for all 
$\lambda>\Lambda:=\max\{1, \Lambda_1, \Lambda_2\}$. This completes the proof.
\end{proof}

Let $\mathcal{J}_{b,\lambda}(u)|_{E_0}$ be a restriction of $\mathcal{J}_{b,\lambda}$ 
on $E_0$, that is
\begin{equation}\label{lab3.23}
\begin{split}
\mathcal{J}_{b,\lambda}(u)|_{E_0}
&=\frac{a}{p}\int_\mathcal{Q}
 \frac{|u(x)-u(y)|^p}{|x-y|^{N+sp}}\,dx\,dy
 +\frac{b}{2p}\Big(\int_\mathcal{Q}\frac{|u(x)-u(y)|^p}{|x-y|^{N+sp}}\,dx\,dy\Big)^2 \\
&\quad -\int_\Omega F(x, u)dx-\int_\Omega G(x, u)dx.
\end{split}
\end{equation}
We define
\begin{align*}
 \widehat{c}=\inf_{\gamma\in\widehat{\Gamma}}
 \max_{t\in[0,1]}  \mathcal{J}_{b,\lambda}|_{E_0} (\gamma(t)),
\end{align*}
where
\[
\widehat{\Gamma}=\{\gamma\in C([0,1],E_0):\gamma(0)=0, \gamma(1)=e\}.
\]
Indeed, it is easily seen that $\widehat{c}$ is independent of $\lambda$. 
Furthermore, if (A4)--(A8) hold, then similarly to the discussion of
 Lemmas \ref{L3.1}--\ref{L3.3} and Remark \ref{rmk3.4}, there exist constants 
$\widehat{b}, \Pi_3>0$ such that for all $b\in(0, \widehat{b})$ and 
$0<|h|_{\frac{p}{p-r}}<\Pi_3$, using Theorem \ref{thm2.4} and 
$E_0\subset X_\lambda$ for all $\lambda>0$, we can conclude that 
$0<\eta\leq c_{b,\lambda}^\beta\leq \widehat{c}$ for all $\lambda\geq\Lambda_1$. 
Now, we can choose $\widehat{M}>\widehat{c}$. Thus
\begin{equation}\label{lab3.24}
0<\eta\leq c_{b,\lambda}^\beta\leq \widehat{c}<\widehat{M}\quad
 \text{for all}\ \lambda\geq\Lambda_1.
\end{equation}


\begin{proof}[Proof of Theorem \ref{thm1.1}]
Let $\beta$ be defined as in Lemma \ref{L3.3}. By Lemmas \ref{L3.1}--\ref{L3.3} 
and Remark \ref{rmk3.4}, there exist constants 
$\Pi:=\min\{\Pi_1,\Pi_2,\Pi_3\}, \widehat{b}>0$ such that for all $b\in(0, \widehat{b})$ 
and $0<|h|_{\frac{p}{p-r}}<\Pi$, $\mathcal{J}_{b,\lambda}$ possesses a $(Ce)$ sequence 
$\{u_n\}\subset X_\lambda$ at the
mountain pass level $c_{b,\lambda}^\beta$ for all $\lambda\geq\Lambda_1$, 
after passing to a subsequence, $\{u_n\}$ satisfying
$$
\sup_{n\in\mathbb{N}^+}\|u_n\|_\lambda\leq\beta,\quad
 \mathcal{J}_{b,\lambda}(u_n)\to  c_{b,\lambda}^\beta\quad \text{and}\quad
 (1+\|u_n\|_\lambda)\|\mathcal{J}_{b,\lambda}'(u_n)\|_{X_\lambda^*}\to  0.
$$
So by Lemma \ref{L3.7}, we can see that there exists $\Lambda>0$ such that 
$\{u_n\}$ has a convergent subsequence in $X_\lambda$ for all 
$\lambda>\Lambda$ and $b\in(0, \widehat{b})$. We may assume that 
$u_n\to  u_{b,\lambda}^{(1)}$ as $n\to \infty$, and thus
$$
0<\|u_{b,\lambda}^{(1)}\|_\lambda\leq\beta,\quad
\mathcal{J}_{b,\lambda}(u_{b,\lambda}^{(1)})= c_{b,\lambda}^\beta\quad \text{and}\quad
\mathcal{J}_{b,\lambda}'(u_{b,\lambda}^{(1)})=0.
$$
Consequently, $u_{b,\lambda}^{(1)}$ is a nontrivial solution for problem \eqref{lab1.1}.

The second nontrivial solution of problem \eqref{lab1.1} will be constructed by the 
local minimization. Now we show that there exists $\psi \in X_\lambda$ such that 
$\mathcal{J}_{b,\lambda}(l \psi)<0$ for all $l>0$ small enough. By \eqref{lab3.1} 
and \eqref{lab3.23}, take $\psi\in E_0$ with $\int_\Omega |\psi|^\nu dx>0$, 
we obtain 
\begin{equation}\label{lab3.25}
\begin{split}
 \mathcal{J}_{b,\lambda}(l\psi)
&= \frac{al^p}{p}\int_\mathcal{Q}\frac{|\psi(x)-\psi(y)|^p}{|x-y|^{N+sp}}\,dx\,dy
+\frac{bl^{2p}}{2p}\Big(\int_\mathcal{Q}\frac{|\psi(x)-\psi(y)|^p}{|x-y|^{N+sp}}\,dx\,dy\Big)^2 \\
&\quad -\int_\Omega F(x, l\psi)dx-\int_\Omega G(x, l\psi)dx \\
&\leq  \frac{al^p}{p}\int_{\mathcal{Q}}\frac{|\psi(x)-\psi(y)|^p}{|x-y|^{N+sp}}\,dx\,dy
+\frac{bl^{2p}}{2p}\Big(\int_\mathcal{Q}\frac{|\psi(x)-\psi(y)|^p}{|x-y|^{N+sp}}\,dx\,dy
 \Big)^2 \\
&\quad -C_1l^\mu\int_{\Omega}|\psi|^\mu dx+C_2l^p\int_{\Omega}|\psi|^pdx-C_3l^\nu\int_{\Omega}|\psi|^\nu dx+C_4|\Omega|  \\
&<0\quad \text{as } l\to 0.
\end{split}
\end{equation}
Then we can deduce that the minimum of the (weakly lower semi-continuous) functional 
$\mathcal{J}_{b,\lambda}$ on any closed ball in $X_\lambda$ with center 0 and radius 
$\varrho<\rho$ satisfying $\mathcal{J}_{b,\lambda}(u)\geq0$ for any $u\in X_\lambda$ 
with $\|u\|_\lambda=\varrho$ is achieved in the corresponding open ball and thus draw 
out a nontrivial solution $u^{(2)}_{b,\lambda}$ of problem \eqref{lab1.1} satisfying 
$\mathcal{J}_{b,\lambda}(u^{(2)}_{b,\lambda})<0$ and $\|u^{(2)}_{b,\lambda}\|<\varrho$. 
In addition, \eqref{lab3.25} means that there exist $l_0>0$ and $\alpha<0$ 
(independent of $\lambda$) such that $\mathcal{J}_{b,\lambda}(l_0 \psi)=\alpha$ and 
$\|l_0 \psi\|<\varrho$. Then it is clear that
\begin{equation}\label{lab3.26}
 \mathcal{J}_{b,\lambda}(u^{(2)}_{b,\lambda})
 \leq\alpha<0<\eta\leq c_{b,\lambda}^\beta=\mathcal{J}_{b,\lambda}(u^{(1)}_{b,\lambda})
\end{equation}
for all $b\in(0, \widehat{b})$, $\lambda>\Lambda$ and $0<|h|_{\frac{p}{p-r}}<\Pi$. 
The proof is complete.
\end{proof}

\begin{thebibliography}{99}
\bibitem[1]{r1} Arosio, A.; Panizzi, S.; \emph{On the well-posedness of the Kirchhoff string}, Trans. Am. Math. Soc., 348 (1996), 305--330.
\bibitem[2]{r2} Bartsch, T.; Wang, Z. Q.; \emph{Existence and multiplicity results for superlinear elliptic problems on R$^N$}, Commun. Partial Differ. Equations, 20 (1995), 1725--1741.
\bibitem[3]{r3} Br\'{e}zis, H.; Lieb, E.; \emph{A relation between pointwise convergence of functions and convergence of functionals}, Proc. Amer. Math. Soc., 88 (1983), 486--490.
\bibitem[4]{r4} Chipot, M.; Lovat, B.; \emph{Some remarks on non local elliptic and parabolic problems}, Nonlinear Anal., 30 (1997), 4619--4627.
\bibitem[5]{r5} Di Nezza, E.; Palatucci, G.; Valdinoci, E.; \emph{Hitchhiker's guide to the fractional Sobolev spaces}, Bull. Sci. math., 136 (2012), 521--573.
\bibitem[6]{r6} Ding, Y. H.; \emph{Variational Methods for Strongly Indefinite Problems}, Interdisciplinary Mathematical Sciences, 7. World Scientific Publishing Co. Pte. Ltd., Hackensack, NJ, 2007.
\bibitem[7]{r7} Ekeland, I.; \emph{Convexity Methods in Hamiltonian Mechanics}, Springer, Berlin, 1990.
\bibitem[8]{r8} Fiscella, A.; Valdinoci, E.; \emph{A critical Kirchhoff type problem involving a nonlocal operator}, Nonlinear Anal., 94 (2014), 156--170.
\bibitem[9]{r9} Goyal, S.; Sreenadh, K.; \emph{Nehari manifold for non-local elliptic operator with concave-convex nonlinearities and sign-changing weight functions}, Proc. Indian Acad. Sci. (Math. Sci.), 125 (2015), 545--558.
\bibitem[10]{r10} He, S. W.; Wen, X. B.; \emph{Existence and concentration of solutions for a Kirchhoff-type problem with sublinear perturbation and steep potential well}, AIMS Math., 8 (2023), 6432--6446.
\bibitem[11]{r11} Iannizzotto, A.; Mosconi, S. J. N.; Squassina, M.; \emph{A note on global regularity for the weak solutions of fractional $p$-Laplacian equations}, Atti Accad. Naz. Lincei Cl. Sci. Fis. Mat. Natur., 27 (2016), 15--24.
\bibitem[12]{r12} Iannizzotto, A.; Mosconi, S. J. N.; Squassina, M.; \emph{Global H\"{o}lder regularity for the fractional $p$-Laplacian}, Rev. Mat. Iberoam., 32 (2016), 1353--1392.
\bibitem[13]{r13} Kirchhoff, G. R.; \emph{Vorlesungen \"{u}ber Mechanik}, Leipzig: B. G. Teubner, 1883.
\bibitem[14]{r14} Laskin, N.; \emph{Fractional quantum mechanics and L\'{e}vy path integrals}, Phys. Lett. A, 268 (2000), 298--305.
\bibitem[15]{r15} Li, Q.; Yang, Z. D.; \emph{Existence and multiplicity of solutions for perturbed fractional $p$-Laplacian equations with critical nonlinearity in $\mathbb{R}^N$}, Appl. Anal., 102 (2023), 2960--2977.
\bibitem[16]{r16} Li, Y. H.; Li, F. Y.; Shi, J. P.; \emph{Existence of a positive solution to Kirchhoff type problems without compactness conditions}, J. Differ. Equations, 253 (2012), 2285--2294.
\bibitem[17]{r17} Liu, S. L.; Chen, H. B.; Yang, J.; Su, Y.; \emph{Existence and nonexistence of solutions for a class of Kirchhoff type equation involving fractional $p$-Laplacian}, RACSAM 114, 161 (2020), 1--28.
\bibitem[18]{r18} Molica Bisci, G.; R\v{a}dulescu, V. D.; Servadei, R.; \emph{Variational Methods for Nonlocal Fractional Problems}, Cambridge Univ. Press, 2016.
\bibitem[19]{r19} Pucci, P.; Xiang, M. Q.; Zhang, B. L.; \emph{Multiple solutions for nonhomogeneous Schr\"{o}dinger-Kirchhoff type equations involving the fractional $p$-Laplacian in $\mathbb{R}^N$}, Calc. Var. Partial Differ. Equations, 54 (2015), 2785--2806.
\bibitem[20]{r20} Pucci, P.; Xiang, M. Q.; Zhang, B. L.; \emph{Existence and multiplicity of entire solutions for fractional $p$-Kirchhoff equations}, Adv. Nonlinear Anal., 5 (2016), 27--55.
\bibitem[21]{r21} Shao, L. Y.; Chen, H. B.; \emph{Multiplicity and concentration of nontrivial solutions for a class of fractional Kirchhoff equations with steep potential well}, Math. Methods Appl. Sci., 45 (2022), 2349--2363.
\bibitem[22]{r22} Servadei, R.; \emph{Valdinoci, E.; Mountain pass solutions for non-local elliptic operators}, J. Math. Anal. Appl., 389 (2012), 887--898.
\bibitem[23]{r23} Sun, J. T.; Wu, T-F.; \emph{Ground state solutions for an indefinite Kirchhoff type problem with steep potential well}, J. Differ. Equations, 256 (2014), 1771--1792.
\bibitem[24]{r24} Szulkin, A.; Weth, T.; \emph{The method of Nehari manifold, Handbook of Nonconvex Analysis and Applications}, DY Gao and D. Motreanu eds, (2010), 597--632.
\bibitem[25]{r25} Tao, H.; Li, L.; Winkert, P.; \emph{Existence and multiplicity of solutions for fractional Schr\"{o}dinger-$p$-Kirchhoff equations in $\mathbb{R}^N$}, Forum Math., 37 (2025), 373--398.
\bibitem[26]{r26} Torres Ledesma, C. E.; \emph{Multiplicity result for non-homogeneous fractional Schr\"{o}dinger-Kirchhoff-type equations in $\mathbb{R}^n$}, Adv. Nonlinear Anal., 7 (2018), 247--257.
\bibitem[27]{r27} Wang, J.; Tian, L. X.; Xu, J. X.; Zhang, F. B.; \emph{Multiplicity and concentration of positive solutions for a Kirchhoff type problem with critical growth}, J. Differ. Equations, 253(2012), 2314--2351.
\bibitem[28]{r28} Wang, L.; Zhang, B. L.; \emph{Infinitely many solutions for Schr\"{o}dinger-Kirchhoff type equations involving the fractional $p$-Laplacian and critical exponent}, Electron. J. Differ. Equations, 339 (2016), 1--18.
\bibitem[30]{r30} Wu, X.; \emph{Existence of nontrivial solutions and high energy solutions for Schr\"{o}dinger-Kirchhoff-type equations in R$^N$}, Nonlinear Anal., 12 (2011), 1278--1287.
\bibitem[31]{r31} Xiang, M. Q.; Zhang, B. L.; Ferrara, M.; \emph{Existence of solutions for Kirchhoff type problem involving the non-local fractional $p$-Laplacian}, J. Math. Anal. Appl., 424 (2015), 1021--1041.
\bibitem[32]{r32} Xiang, M. Q.; Zhang, B. L.; R\u{a}dulescu, V. D.; \emph{Superlinear Schr\"{o}dinger-Kirchhoff type problems involving the fractional $p$-Laplacian and critical exponent}, Adv. Nonlinear Anal., 9 (2019), 690--709.
\bibitem[33]{r33} Xiong, C. W.; Chen, C. F.; Chen, J. H.; Sun, J. J.; \emph{A concave-convex Kirchhoff type elliptic equation involving the fractional $p$-Laplacian and steep well potential}, Complex Var. Elliptic Equations, 68(2023), 932--962.
\bibitem[34]{r34} Xu, L. P.; Chen, H. B.; \emph{Nontrivial solutions for Kirchhoff-type problems with a parameter}, J. Math. Anal. Appl., 433 (2016), 455--472.
\bibitem[35]{r35} Zhang, B. L.; Fiscella, A.; Liang, S. H.; \emph{Infinitely many solutions for critical degenerate Kirchhoff type equations involving the fractional $p$-Laplacian}, Appl. Math. Optim., 80 (2019), 63--80.
\end{thebibliography}
\end{document}
